\documentclass[reqno,11pt]{amsart}
\usepackage{math327-lecture}
\usepackage{needspace}
\newcommand{\Q}{\mathbb{Q}}
\newcommand{\C}{\mathbb{C}}
\newcommand{\F}{\mathbb{F}}
\lectureinfo{26}{Nov. 18, 2026}{Ring Homomorphisms, Substitution, and Ideals}

% Typeset from handwritten Lecture 28, Nov. 19, 2025, pp. 1--12.
\begin{document}
\raggedbottom
\makelecturetitle

\section{Units and fields}

Throughout, $R=(R,+,\cdot,0,1)$ denotes a ring; as in the previous
lecture, rings are assumed commutative with identity.

\begin{define}[Units and fields]
The \emph{group of units} of $R$ is
\[
  R^\times=\{a\in R:\ ab=1\text{ for some }b\in R\}.
\]
If $R^\times=R\setminus\{0\}$, then $R$ is called a \emph{field}.
\end{define}

\begin{example}
The rings $\Z$, $\R$, $\C$, $\Z/n\Z$, and $\F_p=\Z/p\Z$ for a prime $p$
are familiar examples. Among these, $\R$, $\C$, and $\F_p$ are fields,
while $\Z$ is not, since $2$ has no multiplicative inverse in $\Z$.

To see that $\Z/p\Z$ is a field, let $a\not\equiv0\pmod p$. Then
$\gcd(a,p)=1$, so there are integers $m,n$ with
\[
  ma+np=1.
\]
Reducing modulo $p$ gives $ma\equiv1\pmod p$; that is, $a$ has a
multiplicative inverse modulo $p$.

On the other hand, in general $(\Z/n\Z)^\times\neq\Z/n\Z\setminus\{0\}$.
For instance, in $\Z/6\Z$ the class of $2$ has no inverse. Indeed, if
$a\cdot2\equiv1\pmod6$, then multiplying by $3$ gives
\[
  a\cdot(2\cdot3)\equiv3\pmod 6.
\]
Since $2\cdot3=6\equiv0$, this says $3\equiv0\pmod6$, a contradiction.
\end{example}

\begin{example}[Polynomial rings]
If $R$ is a ring, then the set of polynomials with coefficients in $R$
in the variable $x$,
\[
  R[x]=\{f=a_0+a_1x+\cdots+a_nx^n:\ a_i\in R,\ n\geq0\},
\]
is a ring under the usual addition and multiplication of polynomials.
Later in this lecture we focus on the case $F[x]$ with $F$ a field.
\end{example}

\section{Ring homomorphisms}

\begin{define}[Ring homomorphism]
Let $R$ and $S$ be rings. A map $\varphi:R\to S$ is a
\emph{homomorphism} if, for all $a,b\in R$,
\[
  \varphi(a+b)=\varphi(a)+\varphi(b),\qquad
  \varphi(ab)=\varphi(a)\varphi(b),
\]
and $\varphi(1_R)=1_S$.
\end{define}

\begin{example}[Homomorphisms out of $\Z$]
For any ring $R$ there is a unique homomorphism $\Z\to R$. It must send
$1\mapsto1$, and therefore for $n\geq0$
\[
  n\longmapsto\underbrace{1+1+\cdots+1}_{n},
\]
while $-n\mapsto-(1+\cdots+1)$ because a homomorphism preserves
additive inverses.
\end{example}

\section{The substitution principle}

Let $R$ and $S$ be rings and let $\varphi:R\to S$ be a homomorphism.
We ask: is there a homomorphism $\Phi:R[x]\to S$ with
\[
  \Phi(r)=\varphi(r)\quad\text{for all }r\in R,
\]
so that $\Phi$ \emph{extends} $\varphi$? For example, can the
inclusion $\Q\to\C$ be extended to a homomorphism $\Q[x]\to\C$?

Suppose such a $\Phi$ exists, and let
$f=a_0+a_1x+\cdots+a_nx^n\in R[x]$. Then
\begin{align*}
  \Phi(f)&=\Phi(a_0)+\Phi(a_1x)+\cdots+\Phi(a_nx^n)\\
  &=\Phi(a_0)+\Phi(a_1)\Phi(x)+\Phi(a_2)\Phi(x)^2+\cdots
     +\Phi(a_n)\Phi(x)^n\\
  &=\varphi(a_0)+\varphi(a_1)\Phi(x)+\cdots+\varphi(a_n)\Phi(x)^n.
\end{align*}
So $\Phi(f)$, if it exists, is completely determined once we specify
the single element $\Phi(x)\in S$.

\begin{proposition}[Substitution principle]\label{prop:substitution}
Let $\varphi:R\to S$ be a homomorphism. Every homomorphism
$\Phi:R[x]\to S$ extending $\varphi$ is determined by
$\Phi(x)=\alpha\in S$. Moreover, every choice of $\alpha\in S$ gives
such a homomorphism, namely
\[
  \Phi\Bigl(\sum_i a_ix^i\Bigr)=\sum_i\varphi(a_i)\alpha^i.
\]
\end{proposition}
\begin{proof}
Uniqueness was shown above. It remains to check that, for $\alpha\in S$,
the map $\Phi$ defined by the displayed formula is actually a
homomorphism, i.e.\ that
\[
  \Phi(f+g)=\Phi(f)+\Phi(g),\qquad \Phi(1)=1,\qquad
  \Phi(fg)=\Phi(f)\Phi(g).
\]
Additivity follows directly from the additivity of $\varphi$, since
polynomials are added coefficientwise. Also
$\Phi(1)=\varphi(1)=1$.

For multiplicativity, write
\[
  f=a_0+a_1x+\cdots+a_nx^n,\qquad g=b_0+b_1x+\cdots+b_mx^m,
\]
with $a_i,b_j\in R$. Then
\begin{align*}
  \Phi(fg)
  &=\Phi\Bigl(\sum_k\Bigl(\sum_{i+j=k}a_ib_j\Bigr)x^k\Bigr)
   =\sum_k\Bigl(\sum_{i+j=k}\varphi(a_ib_j)\Bigr)\alpha^k\\
  &=\sum_k\Bigl(\sum_{i+j=k}\varphi(a_i)\varphi(b_j)\Bigr)\alpha^k
   =\Bigl(\sum_i\varphi(a_i)\alpha^i\Bigr)
    \Bigl(\sum_j\varphi(b_j)\alpha^j\Bigr)
   =\Phi(f)\,\Phi(g),
\end{align*}
where the second-to-last equality is the expansion of a product of
two sums in the commutative ring $S$.
\end{proof}

\begin{example}[{Evaluation at $\sqrt[3]{2}$}]
Let $\varphi:\Q\to\C$ be the inclusion $\varphi(r)=r$ (usually we do
not give it a name). A homomorphism $\Phi:\Q[x]\to\C$ extending it is
determined by where it sends $x$. For instance, send
\[
  x\longmapsto\gamma:=\sqrt[3]{2},
\]
the real cube root of $2$. Then
\[
  \Phi(x^3-2)=\gamma^3-2=0,\qquad \Phi(x-2)=\gamma-2.
\]
\end{example}

\section{Kernels and ideals}

\begin{define}[Kernel]
If $\varphi:R\to S$ is a homomorphism, its \emph{kernel} is
\[
  \ker(\varphi)=\{x\in R:\ \varphi(x)=0\}.
\]
\end{define}

\begin{warningbox}
In general $\ker(\varphi)$ is \emph{not} a ring: it need not contain
$1$. For example, let $\varphi:\Z\to\Z/p\Z$ be reduction modulo $p$,
$\varphi(x)=x\bmod p$. Then
\[
  \ker(\varphi)=\{0,\pm p,\pm2p,\dots\}=p\Z,
\]
which does not contain $1$.
\end{warningbox}

The kernel is certainly a group under addition (though not a ring),
but it has one more interesting property.

\begin{define}[Ideal]
An \emph{ideal} $I\subseteq R$ is a subset such that
\begin{enumerate}[label=(\arabic*)]
  \item $0\in I$;
  \item $I$ is closed under addition;
  \item if $a\in R$ and $b\in I$, then $ab\in I$.
\end{enumerate}
\end{define}

Conditions (1) and (2), together with (3) applied to $a=-1$, say that
$I$ is an additive subgroup of $R$.

\begin{lemma}\label{lem:kernel-ideal}
If $\varphi:R\to S$ is a homomorphism, then $\ker(\varphi)$ is an ideal
of $R$.
\end{lemma}
\begin{proof}
First, $0\in\ker(\varphi)$ since $\varphi(0)=0$. If
$a,b\in\ker(\varphi)$, then
\[
  \varphi(a+b)=\varphi(a)+\varphi(b)=0+0=0,
\]
so $a+b\in\ker(\varphi)$. Finally, if $a\in\ker(\varphi)$ and
$b\in R$, then
\[
  \varphi(ab)=\varphi(a)\varphi(b)=0\cdot\varphi(b)=0,
\]
so $ab\in\ker(\varphi)$.
\end{proof}

\section{Examples of ideals}

\begin{example}[Ideals of a field]
Let $F$ be a field. Two ideals are always present:
\[
  (0)=\{0\}\qquad\text{and}\qquad (1)=F,
\]
the entire field. These are the only ones. Let $I\subseteq F$ be an
ideal with $I\neq(0)$, and choose $a\in I$ with $a\neq0$. By (3), all
multiples $ba$ with $b\in F$ lie in $I$. In particular,
\[
  a^{-1}a=1\in I.
\]
But if $1\in I$, then every $c\in F$ is also in $I$, since
$c=c\cdot1$. Hence $I=F$.
\end{example}

\begin{exercise}
Check that a ring $R$ with exactly two ideals, $(0)$ and $R$, is a
field.
\end{exercise}

\begin{example}[Ideals of $\Z$]
For $n\geq0$, the sets
\[
  I=n\Z=\{0,\pm n,\pm2n,\dots\}
\]
are clearly ideals of $\Z$. These are the only possibilities: ideals
are in particular subgroups of $(\Z,+)$, and we classified those
subgroups before; each is of the form $n\Z$.
\end{example}

\begin{example}[{Ideals of $\Q[x]$}]
Every ideal in $\Q[x]$ is of the form
\[
  (f)=\{gf:\ g\in\Q[x]\}
\]
for some $f\in\Q[x]$.
% NOTE: The notes end with the fragment "(f_1, f_2, ..., f_n)"; we
% record it as the ideal generated by several elements.
More generally, for $f_1,\dots,f_n\in\Q[x]$ one writes
$(f_1,f_2,\dots,f_n)$ for the ideal they generate.
\end{example}

\end{document}
